\documentclass[conference]{IEEEtran}

\usepackage{amsmath}
\usepackage{amssymb}
\usepackage{bm}
\usepackage{booktabs}
\usepackage{array}
\usepackage{graphicx}
\usepackage{multirow}
\usepackage{url}
\usepackage{hyperref}
\usepackage{xcolor}
\usepackage{listings}
\usepackage{cite}
\usepackage{microtype}
\usepackage{tikz}
\usetikzlibrary{shapes,arrows,positioning,calc}

% Consistent aliases and macros as recommended for ML scientific papers
\newcommand{\tsxtract}{\texttt{Tsxtract}}
\newcommand{\tsxtractor}{\texttt{tsxtractor}}
\newcommand{\tsfresh}{\texttt{tsfresh}}
\newcommand{\catchtwentytwo}{\texttt{catch22}}
\newcommand{\tsfel}{\texttt{TSFEL}}
\newcommand{\hctsa}{\texttt{hctsa}}
\newcommand{\ucrarchive}{UCR Time Series Classification Archive}

% Search paths for graphics: current directory, ./figures/, figures/, paper/figures/, and ./paper/figures/
\graphicspath{{./}{./figures/}{figures/}{paper/figures/}{./paper/figures/}}

% Robust figure inclusion macro: checks local root, figures/, paper/figures/, and formats
\newcommand{\includefigsafe}[3][width=\linewidth]{%
  \IfFileExists{#2}{%
    \includegraphics[#1]{#2}%
  }{%
    \IfFileExists{figures/#2}{%
      \includegraphics[#1]{figures/#2}%
    }{%
      \IfFileExists{paper/figures/#2}{%
        \includegraphics[#1]{paper/figures/#2}%
      }{%
        \IfFileExists{paper/#2}{%
          \includegraphics[#1]{paper/#2}%
        }{%
          \begin{center}
          \begin{tikzpicture}
            \draw[draw=black!40,fill=gray!8,rounded corners=3pt] (0,0) rectangle (7.2,2.6);
            \node[align=center,font=\small\bfseries,text=black!75] at (3.6,1.6) {Figure: #2};
            \node[align=center,font=\footnotesize,text=black!55] at (3.6,1.0) {#3};
            \node[align=center,font=\tiny,text=red!75!black] at (3.6,0.5) {(To display: upload \texttt{#2} or \texttt{figures/} to Overleaf)};
          \end{tikzpicture}
          \end{center}
        }%
      }%
    }%
  }%
}

% Define syntax highlighting for Rust in listings
\lstdefinelanguage{Rust}{
  keywords={as, break, const, continue, crate, else, enum, extern, false, fn, for, if, impl, in, let, loop, match, mod, move, mut, pub, ref, return, self, Self, static, struct, super, trait, true, type, unsafe, use, where, while, async, await, dyn},
  keywordstyle=\color{blue}\bfseries,
  ndkeywords={bool, u8, u16, u32, u64, u128, usize, i8, i16, i32, i64, i128, isize, f32, f64, char, str, Option, Some, None, Result, Ok, Err, String, Vec},
  ndkeywordstyle=\color{purple}\bfseries,
  sensitive=true,
  comment=[l]{//},
  morecomment=[s]{/*}{*/},
  commentstyle=\color{green!50!black}\itshape,
  stringstyle=\color{red!70!black}\ttfamily,
  morestring=[b]",
  morestring=[b]'
}

% Define syntax highlighting for bash in listings
\lstdefinelanguage{bash}{
  keywords={git, pip, python, cargo, checkout, install},
  keywordstyle=\color{blue}\bfseries,
  sensitive=true,
  comment=[l]{\#},
  commentstyle=\color{green!50!black}\itshape,
  stringstyle=\color{red!70!black}\ttfamily,
  morestring=[b]",
  morestring=[b]'
}

\lstset{
  basicstyle=\ttfamily\footnotesize,
  breaklines=true,
  frame=single,
  backgroundcolor=\color{gray!10},
  keywordstyle=\color{blue},
  commentstyle=\color{green!50!black}
}

\begin{document}

\title{\tsxtract{}: High-Throughput Batch and Streaming Time-Series Feature Extraction with a Native Rust Computational Core}

\author{
\IEEEauthorblockN{Aamod Kumar}
\IEEEauthorblockA{\textit{Department of Computer Science and Engineering}\\
\textit{Independent Researcher}\\
Uttar Pradesh, India\\
\url{https://github.com/Aamod007/Tsxtract}}
}

\maketitle

\begin{abstract}
Extracting descriptive mathematical features from time series is a foundational primitive across machine learning, industrial Internet of Things (IoT), and signal processing. However, established Python frameworks such as \tsfresh{}, \catchtwentytwo{}, and \tsfel{} encounter severe throughput bottlenecks when processing large batches ($10^4$--$10^6$ series), driven by serial Foreign Function Interface (FFI) dispatch loops, Global Interpreter Lock (GIL) contention, defensive array copying, and high intra-feature redundancy. This paper investigates whether a batch-oriented systems architecture for classical time-series feature extraction can reduce computational latency and memory overhead at scale while preserving numerical precision and downstream predictive utility. We present \tsxtract{} (\tsxtractor{}), an open-source systems library with a multi-threaded Rust computational core. \tsxtract{} implements zero-copy NumPy buffer ingestion (performing no input-buffer copy during ingestion), releases the GIL, parallelizes work across the series dimension via work-stealing thread pools, and condenses time-domain scans into five cache-conscious fused memory sweeps. Furthermore, it introduces an incremental $O(1)$ stateful accumulator engine for running moments and trend covariance, reducing sliding-window extraction complexity from $O(TW)$ to $O(T)$ for $T$ streamed samples over window capacity $W$. On a benchmark of 1,000 series of length 500 using 16 worker threads on a 10-core/16-thread processor, \tsxtract{} achieves a median throughput of 800,256 series/sec (median execution time of 1.25~ms, IQR: 0.26~ms), outperforming single-series looped C and Python libraries by over two orders of magnitude in feature-normalized latency. Empirical evaluations on synthetic and canonical dynamic time-series tasks demonstrate strong predictive discriminability (96.5\%--100.0\% cross-validated accuracy) with sub-millisecond extraction, while collinearity analysis across diverse dynamic regimes reveals that 83.3\% of feature pairs exhibit low pairwise linear correlation ($|r| < 0.70$).
\end{abstract}

\begin{IEEEkeywords}
Time-series analysis, feature engineering, high-performance computing, Rust, Python, PyO3, Rayon, systems optimization.
\end{IEEEkeywords}

\section{Introduction}
\label{sec:introduction}
Feature-based time-series analysis maps arbitrary-length ordered sequences into fixed-dimensional vectors of statistical, structural, and spectral properties~\cite{fulcher2014highly}. These summary representations enable classical tabular machine learning models---such as Gradient Boosted Trees and Random Forests---to perform classification, clustering, and anomaly detection on temporal data without incurring the prohibitive training and inference overheads of deep recurrent or transformer architectures~\cite{dau2019ucr}.

Over the past decade, several influential software frameworks have emerged to standardize time-series feature engineering. Prominent among them are \tsfresh{}~\cite{christ2018tsfresh}, \catchtwentytwo{}~\cite{lubba2019catch22}, \tsfel{}~\cite{barandas2020tsfel}, and \hctsa{}~\cite{fulcher2014highly}. Despite widespread adoption, existing implementations reveal severe systems-level inefficiencies when scaled to massive batch workloads (e.g., millions of IoT sensor traces, financial tick records, or high-frequency biosignals):
\begin{enumerate}
    \item \textbf{Interpreter Boundary and GIL Contention}: While certain toolkits provide compiled C routines (e.g., \catchtwentytwo{}), their Python wrappers invoke the native kernel sequentially per series inside interpreted Python loops. This pattern incurs repeated FFI marshalling overhead and prevents true multi-core parallel scaling due to the CPython Global Interpreter Lock (GIL).
    \item \textbf{Defensive Memory Duplication}: Frameworks frequently require reshaping tabular inputs into long-format relational DataFrames, creating intermediate memory footprints that exceed the raw data size by an order of magnitude.
    \item \textbf{High Algorithmic Redundancy}: Exhaustive libraries compute hundreds to thousands of feature configurations (e.g., up to 1,558 in \tsfresh{} under \texttt{ComprehensiveFCParameters}, 777 in \texttt{EfficientFCParameters}, and $\sim$390 in \tsfel{}). Dimensionality reduction studies have shown that as few as four principal components explain over 90\% of the variance across these extensive banks~\cite{barandas2020tsfel}, consuming substantial compute budgets on collinear metrics.
    \item \textbf{Sliding-Window Recomputation}: In streaming and rolling-window regimes, naive systems slice windows of length $W$ and recompute all features from scratch on every step, yielding an unfavorable $O(TW)$ total time complexity for $T$ incoming samples.
\end{enumerate}

To address these limitations, this paper investigates the following central research question:
\begin{quote}
\textit{Can a batch-oriented systems architecture for classical time-series feature extraction reduce computational cost and memory overhead at scale while maintaining numerical correctness and practical downstream utility?}
\end{quote}

We emphasize that the primary contribution of \tsxtract{} lies in systems engineering, memory architecture, and batch execution strategy rather than the invention of new mathematical feature estimators. We present the design, algorithmic formulation, and empirical evaluation of \tsxtract{}, an open-source systems library with a high-performance Rust core~\cite{matsakis2014rust} exposed to Python via zero-copy PyO3 bindings~\cite{pyo3_rust_python} for NumPy buffers~\cite{harris2020array}. Specifically, this work makes the following contributions:
\begin{enumerate}
    \item \textbf{High-Throughput Batch Architecture}: A zero-copy buffer ingestion and GIL-detached parallel architecture that processes batches of time series concurrently across worker threads, achieving $>800,000$ series/sec on 16 worker threads.
    \item \textbf{Micro-Architectural Algorithmic Optimizations}: Practical implementations of $O(n)$ Quickselect for 10 quantile ranks, real-to-complex FFT, branchless permutation entropy tables, and five fused memory sweeps.
    \item \textbf{Incremental Streaming Engine}: An online rolling-window state accumulator maintaining circular buffer moments and trend covariance, reducing sliding-window extraction complexity from $O(TW)$ to $O(T)$ for $T$ streamed samples over window capacity $W$.
    \item \textbf{Empirical Evaluation and Correctness}: Comprehensive benchmarking against \tsfresh{}, \catchtwentytwo{}, and \tsfel{}, multi-stage architectural ablation experiments, and 98 automated tests validating numerical agreement with NumPy/SciPy within $10^{-9}$ relative tolerance.
\end{enumerate}

\section{Background and Related Work}
\label{sec:related_work}
Time-series feature extraction libraries generally fall along a spectrum balancing feature exhaustiveness against computational efficiency, as summarized in Table~\ref{tab:comparison_landscape}.

\begin{table*}[t]
\centering
\setlength{\tabcolsep}{3.5pt}
\caption{Comparison of Established Time-Series Feature Extraction Systems. Notice that \tsxtract{} combines native batch-parallel execution across series, zero-copy buffer ingestion, and GIL-independent execution in a Python-facing framework, targeting high-throughput batch pipelines.}
\label{tab:comparison_landscape}
\resizebox{\textwidth}{!}{%
\begin{tabular}{@{}lccccc>{\raggedright\arraybackslash}p{4.5cm}@{}}
\toprule
\textbf{System} & \textbf{Language/Core} & \textbf{Feature Count} & \textbf{Batch Parallelism} & \textbf{Memory Model} & \textbf{GIL Released} & \textbf{Target Niche / Architectural Trade-offs} \\
\midrule
\hctsa{}~\cite{fulcher2014highly} & MATLAB & 7,700+ & Distributed (Cluster) & File / Matrix & N/A & Academic feature discovery; highly exhaustive; computationally intensive. \\
\tsfresh{}~\cite{christ2018tsfresh} & Python & 777 (up to 1,558) & Multiprocessing & Long DataFrame & No & Hypothesis-test feature filtering; slow at scale; high memory duplication. \\
\tsfel{}~\cite{barandas2020tsfel} & Python / Numba & 156 ($\sim$390 avail.) & Serial Loop & Array View & Partial & Domain-specific feature sets; high within-set redundancy; serial FFI dispatch. \\
\catchtwentytwo{}~\cite{lubba2019catch22} & C / Python & 22 & Serial Loop & Pointer Pass & No & Literature-selected canonical features; fast per single series; serial Python loop. \\
\textbf{\tsxtract{}} & \textbf{Rust + PyO3} & \textbf{33} & \textbf{Rayon (16 Threads)} & \textbf{Zero-Copy View} & \textbf{Yes} & \textbf{High-throughput batch and streaming feature extraction; curated non-redundant bank.} \\
\bottomrule
\end{tabular}%
}
\end{table*}

\subsection{Exhaustive Toolkits: hctsa and tsfresh}
Fulcher and Jones~\cite{fulcher2014highly} introduced highly comparative time-series analysis by assembling over 7,700 feature extraction operations in MATLAB (\hctsa{}). While foundational for algorithmic discovery, its computational demand restricts its use on large datasets. Christ et al.~\cite{christ2018tsfresh} adapted these principles into the Python ecosystem with \tsfresh{}, integrating automated feature selection based on hypothesis testing. However, \tsfresh{} provides up to 1,558 feature configurations in its \texttt{ComprehensiveFCParameters} set; our comparative benchmarks evaluate its standard \texttt{EfficientFCParameters} configuration containing 777 features. Computing these exhaustive configurations requires seconds to hours per batch, and its mandatory long-format input structure imposes substantial memory overhead.

\subsection{Curated and Domain Libraries: catch22 and TSFEL}
Recognizing the computational expense of \hctsa{}, Lubba et al.~\cite{lubba2019catch22} developed \catchtwentytwo{}, distilling the 7,700 candidate features down to 22 canonical features through mutual-information clustering on the \ucrarchive{}~\cite{dau2019ucr}. \catchtwentytwo{} is implemented in pure C and provides bindings for Python, R, and Julia. However, its Python bindings operate on individual 1D arrays, necessitating a Python-level iteration loop over batches of time series. Barandas et al.~\cite{barandas2020tsfel} organized features across temporal, statistical, and spectral domains in \tsfel{} using Numba compilation, providing $\sim$390 features in total; when evaluated with all domains enabled at default settings, it computes 156 features. Like \catchtwentytwo{}, \tsfel{} lacks native multi-series parallel batch scheduling in its Python interface.

\subsection{Positioning Tsxtract}
\tsxtract{} is deliberately positioned to occupy a distinct systems niche: \textit{high-throughput batch and online extraction of a curated feature bank with comparatively low pairwise linear correlation on the evaluated corpus over massive collections of time series}. It is not intended to replace \tsfresh{} for exploratory hypothesis screening, nor does it aim to provide multi-language bindings like \catchtwentytwo{}. Instead, it focuses on maximizing hardware utilization, memory locality, and throughput in high-volume Python data pipelines. We emphasize that the core contribution is systems architecture, memory safety, and thread parallelism rather than feature innovation.

\section{System Architecture}
\label{sec:architecture}
\tsxtract{} employs a two-tier architecture: an ergonomic Python wrapper layer, and an optimized Rust execution core compiled as a native C-extension module (\texttt{\_core.pyd} / \texttt{\_core.so}). Figure~\ref{fig:architecture} illustrates the end-to-end processing pipeline.

\begin{figure*}[t]
\centering
\includefigsafe[width=0.96\textwidth]{architecture.png}{End-to-End System Architecture Pipeline}
\caption{System Architecture Pipeline. End-to-end processing pipeline of \tsxtract{} partitioned into Python user space (left) and the native Rust computational core (right). Input matrices cross the language boundary once via unowned borrowed views (\texttt{PyReadonlyArray2}), detaching the CPython GIL. Rayon schedules batches across worker threads along the series axis ($N$), executing five cache-resident sweeps alongside specialized sub-algorithms (Quickselect, RealFFT, branchless permutation entropy) before writing directly into preallocated output buffers. Lower panel highlights the incremental $O(1)$ streaming engine (\texttt{StreamingExtractor}).}
\label{fig:architecture}
\end{figure*}

\subsection{Memory Ingestion and Zero-Copy Borrowing}
When a user calls \texttt{extract\_features(X)}, \texttt{X} is inspected at the PyO3 boundary~\cite{pyo3_rust_python} (\texttt{src/ffi.rs}). If \texttt{X} is a 2D NumPy array of type \texttt{float64}, the runtime utilizes \texttt{numpy::PyReadonlyArray2}, obtaining a borrowed, unowned read-only view into the existing memory buffer:
\begin{lstlisting}[language=Rust]
let view = arr.as_array();
let slice = view.as_slice()
    .ok_or(TsxError::NotContiguous { index: None })?;
\end{lstlisting}
This borrowed view requires that the memory layout be C-contiguous. Non-contiguous or strided inputs raise an explicit \texttt{ValueError} rather than triggering a silent, costly defensive copy. This design guarantees $O(1)$ dynamic memory allocation during ingestion, as the pre-existing input buffer (e.g., 400.0~MB for $N=100,000, n=500$) is borrowed directly without memory duplication.

\subsection{GIL Release and Work-Stealing Parallelism}
CPython's Global Interpreter Lock restricts thread execution to a single core during bytecode interpretation. In \tsxtract{}, immediately following structural input validation, the GIL is released across the entire feature extraction region via PyO3's \texttt{Python::allow\_threads} (wrapped as \texttt{py.detach}):
\begin{lstlisting}[language=Rust]
let flat = py.detach(|| {
    let rows: Vec<&[f64]> = slice.chunks_exact(ncols).collect();
    extract::extract_rows(&rows)
});
\end{lstlisting}
Once detached, the Rust runtime dispatches chunks across threads using Rayon's work-stealing scheduler~\cite{matsakis2015rayon,blumofe1999scheduling}. Parallelism is applied strictly across the \textbf{series dimension} ($N$), rather than across individual feature operations within a single series ($L$). Consequently, each worker thread independently extracts all 33 features for an assigned series, completely eliminating synchronization barriers, mutex contention, and cross-thread communication during compute passes.

\subsection{Intermediate Allocations and Thread-Local Buffers}
To prevent allocator contention in multi-threaded execution, \tsxtract{} eliminates dynamic heap allocations within the inner loops. It maintains two thread-local memory pools:
\begin{enumerate}
    \item \texttt{ORDER\_BUF}: A thread-local \texttt{RefCell<Vec<f64>>} used to partition elements during order-statistic selection without allocating per series.
    \item \texttt{WORKSPACE}: Thread-local real-to-complex FFT scratch buffers and planners that cache execution plans per series length.
\end{enumerate}
Consequently, auxiliary intermediate scratchpad memory scales strictly as $O(p \cdot n)$ with the worker thread count $p$ and sequence length $n$ ($\approx 0.1$~MiB across 16 threads at $n=500$), while output memory scales as $O(N \cdot F)$ ($F=33$), and the input matrix $O(N \cdot n)$ remains unowned resident memory without duplication.

\section{Feature Extraction Methodology}
\label{sec:feature_set}
\tsxtract{} extracts a frozen registry of 33 features per series (\texttt{src/features/mod.rs}). To minimize memory bandwidth consumption, the calculations are reorganized into \textbf{five cache-conscious fused memory sweeps} across the data vector $\bm{x} = [x_0, x_1, \dots, x_{n-1}]^T$:

\subsection{Pass 1: Accumulations, Extremes, and Validity Checks}
A single linear traversal evaluates the raw sums, extreme values, and integrity predicates:
\begin{equation}
\label{eq:pass1_sums}
\text{sum} = \sum_{i=0}^{n-1} x_i, \quad \text{energy} = \sum_{i=0}^{n-1} x_i^2,
\end{equation}
\begin{equation}
\label{eq:pass1_extremes}
\min = \min_{0 \le i < n} x_i, \quad \max = \max_{0 \le i < n} x_i.
\end{equation}
This pass concurrently evaluates whether any value is \texttt{NaN} or if the series is strictly constant ($x_i = x_0, \forall i$). If any \texttt{NaN} is detected, the pipeline short-circuits and assigns \texttt{NaN} to all 33 features, conforming to the library's strict value contract. The population mean is computed as $\mu = \text{sum} / n$.

\subsection{Pass 2: Central Moments}
A second pass accumulates the second, third, and fourth central moment sums relative to $\mu$:
\begin{equation}
\label{eq:pass2_central_moments}
m_k = \sum_{i=0}^{n-1} (x_i - \mu)^k, \quad k \in \{2, 3, 4\}.
\end{equation}
The population variance $\sigma^2$ ($\text{ddof}=0$), standard deviation $\sigma$, skewness $S$, and excess kurtosis $K$ are derived via:
\begin{equation}
\label{eq:pass2_var_std}
\sigma^2 = \frac{m_2}{n}, \quad \sigma = \sqrt{\sigma^2},
\end{equation}
\begin{equation}
\label{eq:pass2_skew_kurt}
S = \frac{m_3 / n}{\sigma^3}, \quad K = \frac{m_4 / n}{\sigma^4} - 3.0.
\end{equation}
For constant series where $\sigma = 0$, $S$ and $K$ analytically resolve to \texttt{NaN}.

\subsection{Pass 3: Successive Differences and Complexity-Invariant Distance}
Successive first differences $\Delta x_i = x_{i+1} - x_i$ are traversed to compute the mean absolute change, mean change, and the Complexity-Invariant Distance ($CID\_CE$)~\cite{batista2014cid}:
\begin{equation}
\label{eq:mean_abs_change}
\text{mean\_abs\_change} = \frac{1}{n-1} \sum_{i=0}^{n-2} |\Delta x_i|,
\end{equation}
\begin{equation}
\label{eq:mean_change}
\text{mean\_change} = \frac{x_{n-1} - x_0}{n-1},
\end{equation}
\begin{equation}
\label{eq:cid_ce}
CID\_CE = \sqrt{\sum_{i=0}^{n-2} \left(\frac{x_{i+1} - x_i}{\sigma}\right)^2}.
\end{equation}
Central second differences $\Delta^2 x_i = x_{i+2} - 2x_{i+1} + x_i$ are averaged to derive $\text{mean\_second\_derivative\_central}$.

\subsection{Pass 4: Thresholds, Crossings, and Run Strikes}
In a fourth fused traversal, level crossings and contiguous run lengths are tracked:
\begin{itemize}
    \item \textbf{Zero Crossings}: Counts indices $i$ where $\text{sgn}(x_i) \neq \text{sgn}(x_{i+1})$.
    \item \textbf{Mean Crossings}: Counts occurrences where $(x_i > \mu) \neq (x_{i+1} > \mu)$.
    \item \textbf{Longest Strike Above/Below Mean}: Evaluates the maximum consecutive run length where $x_i > \mu$ (or $x_i < \mu$).
    \item \textbf{Peak Count}: Counts local maxima with support parameter $s=3$, evaluating $x_i > x_{i \pm d}$ for all $d \in \{1, 2, 3\}$.
\end{itemize}

\subsection{Pass 5: Multi-Lag Autocorrelation and Linear Trend}
Autocorrelation across four discrete lags $\mathcal{K} = \{1, 2, 5, 10\}$ is evaluated in a single fused pass:
\begin{equation}
\label{eq:autocorr}
\rho_k = \frac{\sum_{i=0}^{n-k-1} (x_i - \mu)(x_{i+k} - \mu)}{(n - k) \sigma^2}, \quad \forall k \in \mathcal{K}.
\end{equation}
The linear trend slope $\beta$ and goodness of fit $r^2$ against an implicit time index $t_i = i$ are derived from the covariance:
\begin{equation}
\label{eq:trend_cov}
\text{Cov}(t, x) = \frac{1}{n} \sum_{i=0}^{n-1} \left(i - \frac{n-1}{2}\right)(x_i - \mu),
\end{equation}
\begin{equation}
\label{eq:trend_slope_r2}
\beta = \frac{\text{Cov}(t, x)}{\text{Var}(t)}, \quad r^2 = \frac{\text{Cov}(t, x)^2}{\text{Var}(t) \sigma^2},
\end{equation}
where $\text{Var}(t) = (n^2 - 1) / 12$.

\subsection{Order Statistics via Selection}
Quantile levels $Q = \{0.10, 0.25, 0.50, 0.75, 0.90\}$ are evaluated using NumPy's default linear interpolation method. Rather than performing a full array sort ($O(n \log n)$), \tsxtract{} applies Quickselect (\texttt{select\_nth\_unstable\_by})~\cite{hoare1961algorithm} targeting only the ten bracketing rank indices. By recursing on median ranks, selection depth is bounded logarithmically, executing in $O(n)$ time.

\subsection{Permutation Entropy and Spectral Features}
Normalized permutation entropy (order 3, delay 1)~\cite{bandt2002permutation} is computed by classifying consecutive 3-tuples into one of six ordinal permutations. \tsxtract{} eliminates branch mispredictions by mapping the three pairwise comparisons into a 3-bit integer key $(a \le b) \ll 2 \,|\, (b \le c) \ll 1 \,|\, (a \le c)$ that indexes a static lookup table.

Spectral features (dominant frequency, spectral centroid, spectral entropy) are computed using a real-to-complex Fast Fourier Transform (\texttt{realfft})~\cite{frigo2005design}. Because the input signal is purely real, the negative-frequency spectrum is redundant; \texttt{realfft} exploits Hermitian symmetry to compute an $n/2$ complex spectrum, delivering a measured 1.94$\times$ empirical speedup (3.14~$\mu$s $\to$ 1.62~$\mu$s) relative to general complex FFT algorithms.

\subsection{Stateful Incremental Streaming Engine}
For sliding-window telemetry applications, \tsxtract{} provides an online streaming engine (\texttt{StreamingExtractor} in \texttt{src/features/}\allowbreak\texttt{streaming.rs}) that maintains a circular ring buffer of capacity $W$. Generalizing classical single-pass online moment accumulators~\cite{welford1962note}, as new samples arrive via \texttt{push(val)}, running power sums, linear trend covariance, successive first differences, and zero crossings update incrementally in strictly $O(1)$ arithmetic time:
\begin{equation}
\label{eq:stream_power_sum}
S_k \leftarrow S_k + x_{\text{new}}^k - x_{\text{old}}^k, \quad k \in \{1, 2, 3, 4\},
\end{equation}
\begin{equation}
\label{eq:stream_trend}
T \leftarrow T - (S_{1, \text{prev}} - x_{\text{old}}) + (W - 1) x_{\text{new}}.
\end{equation}
\begin{equation}
\label{eq:stream_diff}
\Delta_{\text{abs}} \leftarrow \Delta_{\text{abs}} + |x_{\text{new}} - x_{\text{prev}}| - |x_{\text{second\_old}} - x_{\text{old}}|.
\end{equation}
This recurrence establishes a two-tiered streaming contract:
\begin{itemize}
    \item \textbf{$O(1)$ Streamable Subset (11 features)}: Online statistical moments ($\mu, \sigma^2, \sigma, S, K$), absolute energy, RMS, linear trend slope $\beta$, trend goodness of fit $r^2$, mean absolute change, and zero crossings are updated in strictly $O(1)$ arithmetic time per incoming sample without rescanning the window.
    \item \textbf{Snapshot Evaluation (22 features)}: Non-linear order statistics (quantiles, median, min, max), run strikes, and spectral transforms evaluate the $W$-element circular ring buffer in $O(W)$ and $O(W \log W)$ time upon explicit request (\texttt{compute\_features}), completely avoiding historical re-allocation.
\end{itemize}

\section{Computational Complexity}
\label{sec:complexity}
Table~\ref{tab:feature_inventory} categorizes the complete inventory of 33 features, detailing asymptotic time and auxiliary memory complexity derived directly from their implementations.

\begin{table}[t]
\centering
\caption{Complete \tsxtract{} Feature Inventory and Algorithmic Complexity. All 33 features operate in $O(n)$ or $O(n \log n)$ time with bounded memory.}
\label{tab:feature_inventory}
\resizebox{\columnwidth}{!}{
\begin{tabular}{@{}llcc>{\raggedright\arraybackslash}p{2.8cm}@{}}
\toprule
\textbf{Group} & \textbf{Feature Names} & \textbf{Time} & \textbf{Memory} & \textbf{Algorithmic Notes} \\
\midrule
\multirow{5}{*}{Stats} & \texttt{mean}, \texttt{std}, \texttt{var} & $O(n)$ & $O(1)$ & 2-pass fused central moments \\
 & \texttt{min}, \texttt{max} & $O(n)$ & $O(1)$ & Fused with Pass 1 sum \\
 & \texttt{median}, \texttt{quantile\_10,25,75,90} & $O(n)$ & $O(n)$ & Quickselect on 10 ranks \\
 & \texttt{skewness}, \texttt{kurtosis} & $O(n)$ & $O(1)$ & Central moments 3 \& 4 \\
 & \texttt{abs\_energy}, \texttt{root\_mean\_square} & $O(n)$ & $O(1)$ & Sum of squares accumulation \\
\midrule
\multirow{2}{*}{Change} & \texttt{mean\_abs\_change}, \texttt{mean\_change} & $O(n)$ & $O(1)$ & 1$^{\text{st}}$ order differences \\
 & \texttt{cid\_ce}, \texttt{mean\_2nd\_deriv\_cent} & $O(n)$ & $O(1)$ & Z-normalized squared diffs \\
\midrule
\multirow{2}{*}{Counts} & \texttt{zero\_crossings}, \texttt{mean\_crossings} & $O(n)$ & $O(1)$ & Fused level detection \\
 & \texttt{number\_of\_peaks}, \texttt{longest\_strike\_*} & $O(n)$ & $O(1)$ & Support $s=3$; run lengths \\
\midrule
\multirow{2}{*}{Correlation} & \texttt{autocorr\_lag\_1,2,5,10} & $O(n)$ & $O(1)$ & Fused 4-lag cross products \\
 & \texttt{trend\_slope}, \texttt{trend\_r2} & $O(n)$ & $O(1)$ & Analytical covariance vs $t$ \\
\midrule
Entropy & \texttt{permutation\_entropy} & $O(n)$ & $O(1)$ & Branchless 3-bit lookup table \\
\midrule
Spectral & \texttt{dominant\_freq}, \texttt{centroid}, \texttt{entropy} & $O(n \log n)$ & $O(n)$ & Real-to-complex FFT \\
\bottomrule
\end{tabular}
}
\end{table}

The composite batch time complexity for processing $N$ independent series of length $n$ across $p$ processor worker threads is:
\begin{equation}
\label{eq:batch_complexity}
T_{\text{batch}}(N, n, p) = O\left(\frac{N \cdot n + N \cdot n \log n}{p}\right) = O\left(\frac{N \cdot n \log n}{p}\right),
\end{equation}
where the $N \cdot n$ component reflects the 30 linear-time time-domain features and the $N \cdot n \log n$ component reflects the three real-to-complex FFT spectral transforms. This bound holds asymptotically under balanced work distribution across worker threads ($N \gg p$), with scheduling overhead amortized. Auxiliary intermediate memory scales strictly as $O(p \cdot n)$, bounded entirely by thread-local scratchpads, while output memory scales as $O(N \cdot F)$ ($F=33$) and resident input memory occupies $O(N \cdot n)$.

\section{Experimental Evaluation}
\label{sec:experiments}

\subsection{Experimental Setup and Hardware}
All benchmark experiments were conducted on a dedicated host workstation:
\begin{itemize}
    \item \textbf{Hardware}: 13th Gen Intel Core i7-13620H processor (hybrid architecture with 10 physical cores: 6 Performance cores + 4 Efficient cores, 16 logical execution threads, up to 4.90~GHz), 16~GB LPDDR5 RAM.
    \item \textbf{Operating System}: Microsoft Windows 11 Enterprise (Build 10.0.26200, x86\_64).
    \item \textbf{Software Stack}: Python 3.14.6, Rust 1.93.0, Cargo 1.93.0, Maturin 1.14.1, NumPy 2.5.2, SciPy 1.17.1, pandas 3.0.3, scikit-learn 1.9.0.
    \item \textbf{Baseline Packages}: \texttt{pycatch22} 0.5.0, \texttt{tsfel} 0.2.0, \texttt{tsfresh} 0.21.2.
\end{itemize}
All timing measurements were gathered on this Windows testbed. Cross-platform CI pipelines on GitHub Actions validate functional correctness and numerical parity across Ubuntu Linux and macOS.

\subsection{End-to-End Batch Pipeline Benchmark Results}
To evaluate real-world pipeline throughput, our benchmark harness (\texttt{benches/\allowbreak bench\_libraries.py}) generates a synthetic batch of 1,000 independent Gaussian series ($N = 1000$, length $n = 500$). To avoid single-run operating system scheduling artifacts or optimistic lower bounds, each library is evaluated after a warm-up pass across repeated iterations ($k=10$ runs for \tsxtract{} and \catchtwentytwo{}, $k \ge 2$ for slower frameworks within a timed budget). We report both the \textbf{Median} and \textbf{Interquartile Range (IQR)}, alongside the minimum wall-clock time. Each library is timed end-to-end including mandatory input ingestion (e.g., long-format relational DataFrames for \tsfresh{}, serial Python loops over 1D series for \catchtwentytwo{} and \tsfel{}).

\begin{table*}[t]
\centering
\setlength{\tabcolsep}{4.5pt}
\caption{End-to-End Batch Feature Extraction Benchmark Results (1,000 Series $\times$ 500 Steps, 16 Worker Threads on 10-Core CPU). Runtimes are reported as Median $\pm$ IQR across repeated runs. Notice that \tsxtract{} achieves over 800,000 series/sec throughput with a feature-normalized latency of 0.0379~ms/feature, outperforming serial Python/C loops by orders of magnitude.}
\label{tab:benchmark_results}
\resizebox{\textwidth}{!}{%
\begin{tabular}{@{}lrrrrrr@{}}
\toprule
\textbf{Library} & \textbf{Feats} & \textbf{Median Time (ms) $\downarrow$} & \textbf{IQR (ms)} & \textbf{Throughput (series/s) $\uparrow$} & \textbf{Latency / Feat (ms) $\downarrow$} & \textbf{Speedup vs Baseline $\uparrow$} \\
\midrule
\textbf{\tsxtract{} 0.3.0} & \textbf{33} & \textbf{1.25} & \textbf{0.26} & \textbf{800,256} & \textbf{0.0379} & \textbf{Baseline (1.0$\times$)} \\
\catchtwentytwo{} (\texttt{pycatch22}) & 22 & 1,024.77 & 1.69 & 976 & 46.5805 & 820$\times$ slower \\
\tsfel{} (all domains) & 156 & 7,154.02 & 45.20 & 140 & 45.8591 & 5,725$\times$ slower \\
\tsfresh{} (\texttt{EfficientFC}) & 777 & 17,683.31 & 120.50 & 57 & 22.7584 & 14,151$\times$ slower \\
\bottomrule
\end{tabular}%
}
\end{table*}

As detailed in Table~\ref{tab:benchmark_results} and Figure~\ref{fig:throughput}, \tsxtract{} processes the entire 1,000-series batch in a median time of \textbf{1.25~ms} (IQR: 0.26~ms), sustaining a throughput of \textbf{800,256 series/sec}. In contrast, \catchtwentytwo{} requires 1.02 seconds, \tsfel{} requires 7.15 seconds, and \tsfresh{} requires 17.68 seconds.

\begin{figure}[t]
\centering
\includefigsafe[width=0.88\linewidth]{throughput.png}{Batch Extraction Throughput across Frameworks}
\caption{Batch Extraction Throughput across Frameworks ($N=1,000$ series $\times 500$ steps, log scale). Error bars denote IQR across 10 repeated benchmark iterations. Notice that \tsxtract{} sustains 800,256 series/sec, driven by zero-copy FFI and multi-threaded work-stealing parallelism.}
\label{fig:throughput}
\end{figure}

\subsection{Deconstructing the Speedup: Matched-Feature Microbenchmark}
A critical scientific question is: \textit{How much of this speedup stems from algorithmic optimization versus batch scheduling, FFI amortization, and multi-core parallelism?} Because toolkits compute distinct feature inventories of varying mathematical complexity, end-to-end runtimes alone cannot establish pure algorithmic superiority.

To address this, we developed a dedicated microbenchmark (\texttt{benches/bench\_matched\_features.py}) measuring the exact per-series execution time ($\mu\text{s}$ per series, median $\pm$ IQR over 10 runs) of identical, semantically equivalent mathematical features on the exact same input matrix ($N=1000, n=500$), reported in Table~\ref{tab:matched_features}.

\begin{table*}[t]
\centering
\setlength{\tabcolsep}{4pt}
\caption{Empirical Matched-Feature Microbenchmark Across Frameworks ($N=1000$ Series $\times$ 500 Steps, Reported as Median $\pm$ IQR in $\mu\text{s}$ per Series). Rows 1--6 report isolated microbenchmark timings of specific estimators on identical inputs. The composite row reports unified multi-threaded batch execution across all 33 features ($1.25 \pm 0.26~\mu$s/series matching Table~\ref{tab:benchmark_results}; in isolated single-thread micro-execution without inter-pass amortization, the accumulated kernel sum is $1.95 \pm 0.68~\mu$s/series).}
\label{tab:matched_features}
\resizebox{\textwidth}{!}{%
\begin{tabular}{@{}lrrrr>{\raggedright\arraybackslash}p{4.2cm}@{}}
\toprule
\textbf{Feature Operation} & \textbf{\tsxtract{} (Batch)} & \textbf{\catchtwentytwo{}} & \textbf{\tsfresh{}} & \textbf{NumPy Ref.} & \textbf{Key Implementation in \tsxtract{}} \\
\midrule
Mean \& Variance & \textbf{0.18 $\pm$ 0.05 $\mu$s} & 15.31 $\pm$ 0.94 $\mu$s & 9.11 $\pm$ 0.62 $\mu$s & 10.24 $\pm$ 0.36 $\mu$s & 2-pass fused central moments \\
Quantiles (5 ranks) & \textbf{0.59 $\pm$ 0.12 $\mu$s} & [Sort $O(n \log n)$] & 130.31 $\pm$ 6.67 $\mu$s & 35.18 $\pm$ 1.64 $\mu$s & Quickselect targeting 10 ranks ($O(n)$) \\
Autocorrelation (Lag 1) & \textbf{0.19 $\pm$ 0.04 $\mu$s} & 72.06 $\pm$ 0.50 $\mu$s & 23.55 $\pm$ 1.80 $\mu$s & 13.26 $\pm$ 1.36 $\mu$s & Fused 4-lag cross products \\
CID-CE (Complexity Distance) & \textbf{0.14 $\pm$ 0.03 $\mu$s} & [Not in set] & 14.55 $\pm$ 1.08 $\mu$s & 12.09 $\pm$ 0.92 $\mu$s & Fused 1$^{\text{st}}$ differences in Pass 3 \\
Zero Crossings & \textbf{0.09 $\pm$ 0.02 $\mu$s} & [Not in set] & 3.73 $\pm$ 0.46 $\mu$s & 4.73 $\pm$ 0.55 $\mu$s & Fused level detection in Pass 4 \\
Permutation Entropy ($d=3, \tau=1$) & \textbf{0.21 $\pm$ 0.05 $\mu$s} & [Not in set] & 278.57 $\pm$ 10.51 $\mu$s & [Not in set] & Branchless 3-bit table lookup \\
\midrule
\textbf{All 33 Features (Parallel Batch)} & \textbf{1.25 $\pm$ 0.26 $\mu$s} & \textbf{1,024.77 $\mu$s (22 feats)} & \textbf{17,683.31 $\mu$s (777 feats)} & \textbf{125.40 $\mu$s (10 feats)} & \textbf{5 fused sweeps + RealFFT + Rayon} \\
\bottomrule
\end{tabular}%
}
\end{table*}

The empirical measurements in Table~\ref{tab:matched_features} reveal that \tsxtract{}'s speedup is not monolithic, but decomposes across four distinct layers:
\begin{enumerate}
    \item \textbf{Algorithmic Optimization ($S_{\text{algorithmic}} \approx 2\times$--$4\times$)}: Quickselect ($O(n)$) achieves a 2.1$\times$ to 59$\times$ reduction in quantile extraction latency relative to full-array sorting in NumPy ($35.18~\mu$s) and \tsfresh{} ($130.31~\mu$s). For permutation entropy ($d=3, \tau=1$), \tsxtract{}'s branchless bitwise mapping ($0.21~\mu$s per series) eliminates branching mispredictions, executing over $1,300\times$ faster than \tsfresh{}'s reference implementation (\texttt{permutation\_entropy} in \texttt{tsfresh.feature\_extraction.feature\_calculators}, $278.57~\mu$s per series). Both implementations were evaluated on identical $n=500$ series with identical parameterization ($d=3, \tau=1$); \tsfresh{} constructs rolling window tuples in Python, hashes permutation sequences in an interpreted dictionary, and calculates Shannon entropy via NumPy, whereas \tsxtract{} maps consecutive 3-tuples into a 3-bit bitwise comparison LUT and accumulates directly into stack-resident registers without heap allocation.
    \item \textbf{Native Loop and Pass Fusion ($S_{\text{native\_loop}} \approx 4.1\times$)}: Condensing calculations into five fused sweeps minimizes cache thrashing, executing 4.1$\times$ faster than an equivalent single-threaded Python/NumPy loop over individual feature calls.
    \item \textbf{FFI Marshalling Amortization ($S_{\text{FFI\_amortization}} \approx 4.1\times$)}: In a controlled single-threaded experiment with identical Rust feature kernels, invoking a single-series native function 1,000 times inside a serial Python loop requires 11.27~ms, whereas a single batch FFI dispatch with a native serial C-loop requires 2.75~ms. This directly isolates an empirical $4.10\times$ FFI boundary and dispatch overhead penalty ($S_{\text{FFI\_amortization}} \approx 4.1\times$). In frameworks such as \texttt{pycatch22}, this repeated per-series Python-to-C wrapper invocation constitutes a major fraction of total runtime.
    \item \textbf{Work-Stealing Multi-Threading ($S_{\text{parallelism}} \approx 13.9\times$)}: Distributing series across 16 worker threads yields an empirical 13.92$\times$ parallel speedup (87.0\% parallel efficiency on a 10-core/16-thread processor).
\end{enumerate}
Rather than assuming strict multiplicative independence among layers, this layered speedup decomposition provides an approximate model consistent with the observed 820$\times$ end-to-end acceleration over \catchtwentytwo{}. It demonstrates how micro-algorithmic efficiency ($2\times\text{--}4\times$), native compiled loop fusion ($4.1\times$), FFI call amortization ($4.1\times$), and thread-level parallelism ($13.9\times$) compound to deliver orders-of-magnitude systems-level gains.

\begin{figure*}[t]
\centering
\includefigsafe[width=0.92\textwidth]{scaling.png}{Scalability Across Dataset Sizes and Sequence Lengths}
\caption{Dual Scaling Evaluation: (a) Batch throughput scaling across dataset sizes ($N = 10^1$ to $10^5$, $n=500$); (b) Sequence length scaling ($n = 50$ to $10,000$, $N=100$) and single-series latency crossover ($N=1$). Notice that \tsxtract{} maintains consistent linear scaling with minimal per-call overhead.}
\label{fig:scaling}
\end{figure*}

\section{Scalability and Memory Analysis}
\label{sec:scalability}

\subsection{Dual Scaling: Batch Size ($N$) and Series Length ($n$)}
Figure~\ref{fig:scaling}(a) illustrates execution time across dataset sizes scaling from $N = 10$ to $N = 100,000$ series ($n = 500$). At small $N \le 100$, thread pool dispatch latency constitutes a modest fraction of total runtime. Beyond $N \ge 1,000$, throughput scales near-linearly with series count, maintaining constant per-series latency.

To evaluate how execution time scales with input duration, Figure~\ref{fig:scaling}(b) benchmarks performance across sequence lengths $n \in \{50, 100, 250, 500, 1000, 2500, 5000, 10000\}$ for a batch of $N = 100$ series. Batch execution time scales near-linearly from 0.084~ms at $n=50$ (0.84~$\mu$s/series) to 3.336~ms at $n=10,000$ (33.36~$\mu$s/series), demonstrating that linear-time fused sweeps dominate practical sequence lengths, with spectral $O(n \log n)$ transforms introducing only a modest logarithmic curve at $n \ge 5,000$.

Furthermore, evaluating single-series latency ($N = 1$) isolates the crossover point between \tsxtract{} and standard NumPy calls. For short single series ($n \le 500$), \tsxtract{} computes all 33 features in 1.6~$\mu$s to 21.4~$\mu$s, executing faster than an 8-feature Python/NumPy loop (36.6~$\mu$s to 40.5~$\mu$s). At $n \approx 1,000$, the crossover occurs: single-series latency is 44.3~$\mu$s for \tsxtract{} versus 42.5~$\mu$s for NumPy. At $n \ge 5,000$, single-series latency reaches 127.8~$\mu$s for \tsxtract{} versus 96.4~$\mu$s for NumPy. This confirms that while native batch multi-threading provides massive throughput scaling across batches ($N \ge 100$), single-series latency at large $n$ is governed by the comprehensive 33-feature workload rather than thread parallelism.

\subsection{Disambiguated Memory Footprint and Zero-Copy Benefits}
To provide an accurate accounting of memory consumption and clearly distinguish pre-existing input buffers from dynamic heap allocations, Figure~\ref{fig:memory} explicitly separates three memory tiers:
\begin{enumerate}
    \item \textbf{Resident Input Buffer}: For a batch of $N = 100,000$ series of length $n = 500$, the input array occupies $100,000 \times 500 \times 8\text{ bytes} = 400.0\text{ MB}$ ($381.47\text{ MiB}$) of pre-existing memory. \tsxtract{} borrows an unowned pointer view via PyO3, performing no input-buffer copy during ingestion (\textbf{zero additional memory duplication}).
    \item \textbf{Allocated Output Matrix}: The extracted feature table occupies $100,000 \times 33 \times 8\text{ bytes} = 26.4\text{ MB}$ ($25.18\text{ MiB}$).
    \item \textbf{Intermediate Heap Allocations}: \tsxtract{} requires only $O(p \cdot n) \approx 0.1\text{ MiB}$ across all $p=16$ worker threads for thread-local scratchpads.
\end{enumerate}
In total, \tsxtract{}'s additional allocated memory footprint is strictly \textbf{25.2~MiB} (26.4~MB), in contrast to frameworks like \tsfresh{} which duplicate input data into long-format indexed DataFrames and allocate intermediate object graphs exceeding \textbf{825~MiB}.

\begin{figure}[t]
\centering
\includefigsafe[width=0.85\linewidth]{memory.png}{Intermediate Memory Footprint (Zero-Copy vs. Duplication)}
\caption{Disambiguated Memory Footprint vs Batch Size ($N$). Notice that the pre-existing 400.0~MB input matrix is borrowed zero-copy, so \tsxtract{} allocates only 25.2~MiB (26.4~MB) for the output matrix, avoiding the $>825$~MiB footprints of defensive copying frameworks.}
\label{fig:memory}
\end{figure}

\subsection{Multi-Core Speedup and Efficiency}
Figure~\ref{fig:speedup} evaluates parallel speedup $S(p) = T_1 / T_p$ and parallel efficiency $\eta(p) = S(p) / p$ across thread counts $p \in \{1, 2, 4, 8, 16\}$. Because series extractions are completely decoupled across the batch axis, \tsxtract{} achieves a parallel speedup of 13.92$\times$ on 16 worker threads ($p=16$), reflecting an 87.0\% parallel efficiency.

\begin{figure}[t]
\centering
\includefigsafe[width=\linewidth]{speedup.png}{Multi-Core Speedup ($S(p)$) and Rayon Efficiency ($\eta(p)$)}
\caption{Multi-Core Speedup ($S(p)$) and Rayon Efficiency ($\eta(p)$) across 16 Worker Threads. Error bars denote IQR across 10 repeated runs. Notice that \tsxtract{} maintains 87.0\% parallel efficiency at 16 threads due to complete series-level independence.}
\label{fig:speedup}
\end{figure}

\section{Ablation Study}
\label{sec:ablation}

\subsection{Architectural Systems Ablation}
To quantify the isolated performance contribution of each architectural layer, we conducted a controlled stage-by-stage ablation on $N = 1000$ series of length $n = 500$ across 10 repeated runs:
\begin{enumerate}
    \item \textbf{Stage 1: Python Baseline Loop}: Sequential Python loop calling NumPy C-extensions per series required \textbf{46.31 $\pm$ 0.61~ms} (1.0$\times$ baseline).
    \item \textbf{Stage 2: Serial FFI Dispatch}: Calling the Rust kernel per series inside a Python loop (1,000 FFI crossings) executed in \textbf{11.27 $\pm$ 0.59~ms} (4.11$\times$ speedup), establishing the raw single-thread native kernel advantage.
    \item \textbf{Stage 3: Batch FFI + Rayon Multi-Threading (16 Threads)}: Ingesting the batch in a single FFI crossing, detaching the GIL, and distributing series across 16 worker threads reduced wall time to \textbf{1.52 $\pm$ 0.32~ms} (\textbf{30.51$\times$ speedup} over Python baseline, and 7.42$\times$ faster than serial FFI dispatch).
    \item \textbf{Stage 4: Defensive Array Copy Overhead}: Forcing defensive memory duplication prior to extraction increased wall time to \textbf{2.61 $\pm$ 0.31~ms} (+1.09~ms copy latency, 72\% slower than zero-copy borrowing).
\end{enumerate}

\subsection{Micro-Architectural Algorithmic Ablation}
Within the native execution core, micro-architectural design choices yield substantial empirical benefits:
\begin{itemize}
    \item \textbf{Quickselect vs Full Sort}: Quickselect targeting 10 quantile ranks reduced quantile extraction time from 4.82~$\mu$s to 2.29~$\mu$s (2.1$\times$ improvement) relative to full sorting.
    \item \textbf{RealFFT vs Complex FFT}: Exploiting Hermitian symmetry halved the transform length, reducing spectral stage latency from 3.14~$\mu$s to 1.62~$\mu$s (1.94$\times$ empirical speedup).
    \item \textbf{Branchless Permutation Entropy}: Eliminating branch mispredictions via static bitwise comparison lookup tables reduced permutation entropy latency from 2.11~$\mu$s to 0.67~$\mu$s per series.
    \item \textbf{Memory Pass Fusion}: Condensing 33 feature passes into five fused contiguous sweeps reduced the number of major sequential feature-computation sweeps from 33 conceptual passes down to five fused memory traversals, maximizing cache line reuse while data elements reside in L1/L2 data cache.
\end{itemize}

\section{Downstream Machine Learning Utility}
\label{sec:downstream_ml}
To verify that \tsxtract{} retains high discriminative representation capacity on non-trivially separable tasks, we evaluated classification performance across four synthetic and canonical dynamic time-series classification tasks with per-series z-score standardization:
\begin{enumerate}
    \item \textbf{Synthetic Control (6-class, $N=600$, $n=150$)}: White noise, cyclic, upward trend, downward trend, upward shift, downward shift with overlapping noise.
    \item \textbf{Simulated ECG (2-class, $N=500$, $n=150$)}: Normal sinus rhythm vs arrhythmic waveforms with physiological baseline wander.
    \item \textbf{Kinematic Gesture (2-class, $N=400$, $n=150$)}: Inertial motion capture gesture profiles (GunPoint equivalent).
    \item \textbf{Power Demand (2-class, $N=400$, $n=150$)}: Diurnal weekday vs weekend telemetry load profiles.
\end{enumerate}

Features were extracted using \tsxtract{} (33 features), \catchtwentytwo{} (22 features), and a naive summary statistics baseline (5 features: mean, std, min, max, median). Both \textbf{Random Forest} (100 estimators) and \textbf{Ridge Classifiers} ($\alpha=1.0$) were evaluated using 5-fold stratified cross-validation.

\begin{table*}[t]
\centering
\setlength{\tabcolsep}{5pt}
\caption{Downstream Machine Learning Classification Results (5-Fold Stratified Cross-Validation). Notice that \tsxtract{} achieves strong classification performance across both linear and non-linear models while executing 252$\times$--631$\times$ faster than \catchtwentytwo{}.}
\label{tab:downstream_results}
\resizebox{\textwidth}{!}{%
\begin{tabular}{@{}llcccccc@{}}
\toprule
\textbf{Dataset} & \textbf{Representation} & \textbf{Feats} & \textbf{Extract Time (ms) $\downarrow$} & \textbf{RF Acc (\%) $\uparrow$} & \textbf{RF F1 $\uparrow$} & \textbf{Ridge Acc (\%) $\uparrow$} & \textbf{Ridge F1 $\uparrow$} \\
\midrule
\multirow{3}{*}{Synthetic Control (6-class)} & \textbf{\tsxtract{}} & \textbf{33} & \textbf{0.61} & \textbf{96.5\%} & \textbf{0.965} & \textbf{93.3\%} & \textbf{0.933} \\
 & \catchtwentytwo{} & 22 & 232.24 & 91.0\% & 0.909 & 85.5\% & 0.854 \\
 & Naive Stats & 5 & 1.03 & 93.5\% & 0.935 & 77.0\% & 0.757 \\
\midrule
\multirow{3}{*}{Simulated ECG (2-class)} & \textbf{\tsxtract{}} & \textbf{33} & \textbf{0.42} & \textbf{100.0\%} & \textbf{1.000} & \textbf{100.0\%} & \textbf{1.000} \\
 & \catchtwentytwo{} & 22 & 191.55 & 100.0\% & 1.000 & 100.0\% & 1.000 \\
 & Naive Stats & 5 & 0.74 & 75.2\% & 0.751 & 62.8\% & 0.626 \\
\midrule
\multirow{3}{*}{Kinematic Gesture (2-class)} & \textbf{\tsxtract{}} & \textbf{33} & \textbf{0.24} & \textbf{100.0\%} & \textbf{1.000} & \textbf{100.0\%} & \textbf{1.000} \\
 & \catchtwentytwo{} & 22 & 151.43 & 100.0\% & 1.000 & 100.0\% & 1.000 \\
 & Naive Stats & 5 & 0.52 & 95.2\% & 0.952 & 96.5\% & 0.965 \\
\midrule
\multirow{3}{*}{Power Demand (2-class)} & \textbf{\tsxtract{}} & \textbf{33} & \textbf{0.61} & \textbf{100.0\%} & \textbf{1.000} & \textbf{100.0\%} & \textbf{1.000} \\
 & \catchtwentytwo{} & 22 & 153.58 & 100.0\% & 1.000 & 100.0\% & 1.000 \\
 & Naive Stats & 5 & 0.60 & 64.2\% & 0.641 & 65.2\% & 0.652 \\
\bottomrule
\end{tabular}%
}
\end{table*}

As shown in Table~\ref{tab:downstream_results}, \tsxtract{} consistently matches or outperforms \catchtwentytwo{} while extracting features in under 0.65~ms across all datasets (252$\times$--631$\times$ faster than \catchtwentytwo{}'s 151--232~ms). Crucially, on tasks where simple moment statistics collapse (e.g., Ridge accuracy falling to 62.8\% on ECG and 65.2\% on Power Demand due to shape ambiguity), \tsxtract{}'s temporal and spectral features retain full discriminability (100.0\% accuracy).

\section{Information-Theoretic Redundancy Analysis}
\label{sec:redundancy}
To empirically assess the degree of collinearity within \tsxtract{}'s feature bank, we evaluated pairwise Pearson correlations across a diverse benchmark corpus of 2,000 synthetic time series spanning eight dynamic regimes (white noise, harmonic, random walk, AR(1), linear trend, regime shift, chaotic logistic maps, and sparse pulses).

Across all 528 unique feature pairs, \textbf{83.3\%} exhibit low pairwise linear correlation ($|r| < 0.70$), 11.6\% exhibit moderate correlation ($0.70 \le |r| < 0.90$), and only 5.1\% exhibit high collinearity ($|r| \ge 0.90$). Principal Component Analysis (PCA) demonstrated that capturing 90\% of total variance requires 6 principal components (18.2\% of the feature bank). Importantly, pairwise linear correlation does not rule out non-linear dependencies, nor does PCA variance retention directly dictate downstream classification utility; nonetheless, this empirical profile demonstrates that \tsxtract{} avoids the excessive collinearity of over-parameterized libraries where as few as 4 components explain over 90\% of an entire 390-feature inventory~\cite{barandas2020tsfel}.

\section{Numerical Correctness and Test Suite}
\label{sec:correctness}
Numerical reliability is validated against NumPy and SciPy reference implementations through an automated suite of 98 unit, property, and regression tests:
\begin{itemize}
    \item \textbf{Reference Agreement}: Every feature is verified to match NumPy/SciPy reference calculations within a relative tolerance of $10^{-9}$ across standard, periodic, and constant data.
    \item \textbf{Property-Based Fuzz Testing}: Property tests using \texttt{hypothesis} generate random series with adversarial floating-point extremes (signed zeros, denormals, infinities) to assert that no panics cross the FFI boundary.
    \item \textbf{Strict NaN Contract}: Any input containing \texttt{NaN} propagates to assign \texttt{NaN} to all 33 features. Individually undefined features (e.g., autocorrelation of constant series) resolve to \texttt{NaN} independently without corrupting companion features.
\end{itemize}

\section{Limitations and Threats to Validity}
\label{sec:threats_to_validity}

\subsection{Limitations}
\begin{enumerate}
    \item \textbf{Non-Exhaustive Feature Space}: \tsxtract{} computes 33 features. It cannot serve workflows requiring exhaustive exploratory screenings of thousands of specialized hypothesis tests.
    \item \textbf{Single-Series Regimes}: For a single short series, native multi-core batching offers negligible advantage over standard NumPy routines, as detailed in Section~\ref{sec:scalability}.
    \item \textbf{Memory Contiguity Constraints}: \tsxtract{} requires C-contiguous \texttt{float64} arrays; non-contiguous arrays require prior normalization.
\end{enumerate}

\subsection{Threats to Validity}
\begin{itemize}
    \item \textbf{Construct Validity}: While milliseconds-per-feature provides a metric of computational density, libraries compute estimators of varying complexity. We mitigated this by conducting rigorous matched-feature microbenchmarks on identical mathematical estimators (Table~\ref{tab:matched_features}).
    \item \textbf{Internal Validity}: Timing benchmarks were collected across repeated iterations ($k=10$) following warm-up passes, reporting medians and IQRs to prevent operating system scheduling noise from distorting measurements.
    \item \textbf{External Validity}: Benchmarks were executed on an x86\_64 host under Windows 11; cross-platform CI workflows verify numerical parity and test suite execution on Ubuntu Linux and macOS.
\end{itemize}

\section{Reproducibility}
\label{sec:reproducibility}
To ensure complete reproducibility, all artifacts, scripts, and tests are released under the MIT license at:
\begin{center}
\url{https://github.com/Aamod007/Tsxtract}
\end{center}
Benchmarks execute with random seed 42, preceded by two untimed warmup iterations per framework to mitigate cache cold-start effects, reporting median $\pm$ IQR over $k=10$ measured runs ($k \ge 2$ for slow baselines). Experiments were executed on an Intel Core i7-13620H processor running Windows 11 Enterprise with OS high-performance power plan enabled and non-essential background processes terminated. The benchmark suite and paper figures can be reproduced via:
\begin{lstlisting}[language=bash]
git checkout 1e8b277
pip install -e ".[bench]"
python benches/bench_libraries.py --n-series 1000 --n-steps 500
python benches/bench_ucr_downstream.py
python paper/generate_figures.py
\end{lstlisting}

\section{Conclusion}
\label{sec:conclusion}
This work examined whether a batch-oriented systems architecture for classical time-series feature extraction can overcome the scaling bottlenecks of existing Python frameworks. The results demonstrate that by combining zero-copy FFI buffer ingestion, GIL release, multi-threaded work-stealing parallelism across series, and cache-conscious fused memory passes, \tsxtract{} achieves an extraction throughput exceeding 800,000 series/sec across 16 worker threads on a 10-core CPU. Furthermore, the integration of an $O(1)$ streaming state engine enables real-time sliding-window processing for running moments and trend covariance. By demonstrating that high extraction throughput can be achieved alongside numerical correctness and strong representation capacity on canonical benchmark tasks, \tsxtract{} provides a practical systems foundation for scalable time-series feature engineering.

% Fully self-contained IEEE bibliography: guarantees 0 errors on Overleaf / TeX Live without needing external .bib files
\begin{thebibliography}{15}

\bibitem{fulcher2014highly}
B.~D. Fulcher and N.~S. Jones, ``Highly comparative feature-based time-series classification,'' \emph{IEEE Transactions on Knowledge and Data Engineering}, vol.~26, no.~12, pp. 3026--3037, 2014.

\bibitem{dau2019ucr}
H.~A. Dau, E.~Keogh, K.~Kamgar, C.-C.~M. Yeh, Y.~Zhu, S.~Gharghabi, C.~A. Ratanamahatana, C.~Yanping, H.~Bing, B.~Nurjahan, \emph{et~al.}, ``The {UCR} time series classification archive,'' \emph{IEEE/CAA Journal of Automatica Sinica}, vol.~6, no.~6, pp. 1293--1305, 2019.

\bibitem{christ2018tsfresh}
M.~Christ, N.~Braun, J.~Neuffer, and A.~W. Kempa-Liehr, ``Time series feature extraction on basis of scalable hypothesis tests (tsfresh -- a {Python} package),'' \emph{Neurocomputing}, vol. 307, pp. 72--77, 2018.

\bibitem{lubba2019catch22}
C.~H. Lubba, S.~S. Sethi, P.~Knaute, S.~R. Schultz, B.~D. Fulcher, and N.~S. Jones, ``catch22: {CAnonical Time-series CHaracteristics}: Selected through comprehensive analysis of 7,000 empirical time-series features,'' \emph{Data Mining and Knowledge Discovery}, vol.~33, no.~6, pp. 1821--1852, 2019.

\bibitem{barandas2020tsfel}
M.~Barandas, D.~Folgado, R.~Fernandes, S.~Santos, M.~Abreu, P.~Bota, H.~Liu, T.~Schultz, and H.~Gamboa, ``{TSFEL}: Time series feature extraction library,'' \emph{SoftwareX}, vol.~11, p. 100456, 2020.

\bibitem{batista2014cid}
G.~E. A. P.~A. Batista, E.~J. Keogh, O.~M. Tataw, and V.~M. A. De~Souza, ``{CID}: An efficient complexity-invariant distance for time series,'' \emph{Data Mining and Knowledge Discovery}, vol.~28, no.~3, pp. 634--669, 2014.

\bibitem{bandt2002permutation}
C.~Bandt and B.~Pompe, ``Permutation entropy: A natural complexity measure for time series,'' \emph{Physical Review Letters}, vol.~88, no.~17, p. 174102, 2002.

\bibitem{frigo2005design}
M.~Frigo and S.~G. Johnson, ``The design and implementation of {FFTW3},'' \emph{Proceedings of the IEEE}, vol.~93, no.~2, pp. 216--231, 2005.

\bibitem{hoare1961algorithm}
C.~A.~R. Hoare, ``Algorithm 65: Find,'' \emph{Communications of the ACM}, vol.~4, no.~7, pp. 321--322, 1961.

\bibitem{matsakis2015rayon}
N.~D. Matsakis, ``Rayon: Data parallelism in {Rust},'' 2015. [Online]. Available: \url{https://github.com/rayon-rs/rayon}

\bibitem{blumofe1999scheduling}
R.~D. Blumofe and C.~E. Leiserson, ``Scheduling multithreaded computations by work stealing,'' \emph{Journal of the ACM}, vol.~46, no.~5, pp. 720--748, 1999.

\bibitem{pyo3_rust_python}
{PyO3 Contributors}, ``{PyO3}: Rust bindings for the {Python} interpreter,'' 2024. [Online]. Available: \url{https://github.com/PyO3/pyo3}

\bibitem{matsakis2014rust}
N.~D. Matsakis and F.~S. Klock, ``The {Rust} language,'' \emph{ACM SIGAda Ada Letters}, vol.~34, no.~3, pp. 103--104, 2014.

\bibitem{harris2020array}
C.~R. Harris, K.~J. Millman, S.~J. van~der Walt, R.~Gommers, P.~Virtanen, D.~Cournapeau, E.~Wieser, J.~Taylor, S.~Berg, N.~J. Smith, \emph{et~al.}, ``Array programming with {NumPy},'' \emph{Nature}, vol. 585, no. 7825, pp. 357--362, 2020.

\bibitem{welford1962note}
B.~P. Welford, ``Note on a method for calculating corrected sums of squares and products,'' \emph{Technometrics}, vol.~4, no.~3, pp. 419--420, 1962.

\end{thebibliography}

\end{document}
